\itemhead{第1章 小结与复习 · 第4题（教材p.28）}{XJ-B1-C01-SU-P028-Y1361}
\label{XJ-B1-C01-SU-P028-Y1361}
{\footnotesize 湘教版教材p.28 / PDF第35页 / 第1章 小结与复习 / 第4题}
\field{原题}{举例说明什么叫作命题，如何对一个命题作否定。}
\field{参考答案}{示例：“任意实数的平方非负”是真命题；否定为“存在实数的平方小于0”，是假命题。否定作用于整个断言，不能只替换一句中的个别词。}
{\small \field{关联知识}{L01 命题、真假、逆命题与否定（教材帮p.30 / PDF39）；L04 含一个量词命题的否定（教材帮p.39、40 / PDF48、49）}
\field{方法与易错点}{明确论域与断言；真命题给理由，假命题给满足题设而使结论失败的反例。 易错：把逆命题当成否定；只改变结论中的个别词而改变原逻辑范围。}
\field{答案依据}{答案册未提供本栏目的独立对应答案；依据原题、教材概念补写参考回应。开放性总结不设唯一数值答案。}}
